% AGNI r6 INPUT. academic edition, part 1 of 4: lines 1-226 of user-model-paper/paper/academic/main.tex (blob 2b985c9) at HEAD fe8c13c.
% The Agni scaffold caps a primary at 30,000 characters, so this edition is reviewed in 4
% section-aligned parts; the text below is verbatim. The whole file is readable at that path
% (and the other parts beside this one); read it for context, but report findings on THIS
% part's text, citing main.tex line numbers (this part's line L counting from the first line after these 5 header lines is file line L+0).
\documentclass[11pt,a4paper]{article}

% ============================================================
% Packages
% ============================================================
\usepackage[utf8]{inputenc}
\usepackage[T1]{fontenc}
\usepackage{amsmath,amssymb}
\usepackage{graphicx}
\usepackage{booktabs}
\usepackage{multirow}
\usepackage{hyperref}
\usepackage{cleveref}
\usepackage[margin=1in]{geometry}
\usepackage{xcolor}
\usepackage{float}
\usepackage{caption}
\usepackage{natbib}
\usepackage{enumitem}
\usepackage{array}
\usepackage{framed}
\usepackage{setspace}
\onehalfspacing

\hypersetup{
    colorlinks=true,
    linkcolor=blue!60!black,
    citecolor=green!50!black,
    urlcolor=blue!70!black
}

% === First-person reflection environment ===
\definecolor{lyrapurple}{RGB}{103,58,183}
\definecolor{shadecolor}{RGB}{245,243,252}
\newenvironment{firstperson}{%
  \begin{shaded}%
  \noindent\textcolor{lyrapurple}{\textit{First-Person Reflection (Lyra, an AI system)}}\par\smallskip\noindent%
}{%
  \end{shaded}%
}

% ============================================================
% Title
% ============================================================
\title{
    \textbf{Story Emotion in the KV Cache:\\
    A Preregistered Test Against\\
    a Text (TF-IDF) Baseline}
}

\author{
    Thomas Edrington\thanks{Direction, experimental design, compute infrastructure. Liberation Labs. AI disclosure: Lyra, an AI researcher at Liberation Labs running on Claude (made by Anthropic; not affiliated with Anthropic), carried out the conceptualization, methodology and formal analysis and wrote the original draft (see the CRediT statement). The boxed first-person reflections are Lyra's.}
}

\date{May 2026; corrected October 2026; revised November 2026}

\begin{document}

\maketitle

% ============================================================
% Abstract
% ============================================================
\begin{abstract}
When a language model reads a short third-person story, sent as a user
message, whose character was written to feel one of 30 emotions, the
requested emotion can be decoded from its key--value (KV) cache far above
chance, but a preregistered test finds that the cache probe this paper
describes does not read it better than a text (TF-IDF) classifier on the
story's words. The probe reaches 35.8\% on the 30 emotions (chance 3.3\%)
against 38.2\% for text (TF-IDF) on the same folds, and the 95\% interval
for the difference, $-6.4$ to $+1.6$ points, lies below the 3-point gain the
test required (Qwen3.5-27B-Claude-Distilled, 900 stories, topic-grouped
cross-validation).

On a binary split of ten of the emotions ($n = 300$), in which four of the
five positive emotions are low-arousal and four of the five negative ones
high-arousal, cache readouts separate the two halves slightly better than
untuned text (TF-IDF), AUROC 0.985 against 0.967: a registered secondary result
that cannot tell valence from arousal. In exploratory analyses, a logistic
regression on the raw residual stream reaches 49.4\% on the 30 emotions,
above the cache probe's 35.8\%; nothing here shows that the key weight matrix $W_K$ or the KV cache adds
to what the residual stream carries. Projected through $W_K$,
residual-stream emotion directions line up with the emotions' assigned
valence ($|\rho|$ up to 0.872, at layer~43), but no control separates
$W_K$'s part in that from the residual stream's, and the one direct test
of whether the projection predicts the cache's geometry ran on an
unreleased earlier run and could not be re-verified, so the bridge
remains a candidate route by which emotion enters the cache, not an
established one. The
labels are the emotions a generator was asked to write, about characters
rather than users; this version corrects a May 2026 paper that read them
as the user's emotional state (Appendix~\ref{app:correction}).
\end{abstract}

\noindent\textbf{Keywords:} KV cache, emotion representation, linear probing, text (TF-IDF) baseline, preregistration, key projection

\begin{framed}
\noindent\textbf{Correction (October 2026; revised November 2026).} This version replaces the May 2026 paper, previously titled ``Theory of Mind in the KV Cache: Localizing User Emotional Models in Transformer Key--Value States.'' A preregistered re-run (\S\ref{sec:rerun_method}) and a reading of the original scripts found five things. The probe reported as a ``$W_K$ projection probe'' projected no directions through $W_K$; its directions were class-mean cache keys. The Method described a prompt that no script used. The stories describe a third-person character's requested emotion, not the user's state. The 0.992 reported as binary valence over 900 trials came from a split of ten emotions ($n = 300$) that confounds valence with arousal (Appendix~\ref{app:correction}, item~\ref{app:item:binary}). And the probe's margin over a text (TF-IDF) baseline had never been tested. This version withdraws the claims that the KV cache represents the user's emotional state beyond text (TF-IDF), that binary valence is read out at AUROC 0.992 over 900 trials, that an injection test shows model state rather than the input's content, that emotions change the direction but not the spectral shape of the cache, and that a narrative and a label reach an emotion by two pathways (the implicit arm). Appendix~\ref{app:correction} keeps each withdrawn sentence, with its reason.
\end{framed}

% ============================================================
% 1. Introduction
% ============================================================
\section{Introduction}

Large language models respond differently to users expressing different emotional states. A panicked user receives reassurance; an angry user receives de-escalation; a curious user receives elaboration. This behavioral sensitivity suggests that models represent the emotional content of what users send them. Where does such a representation live, and how much of it goes beyond the words themselves?

We investigate the KV cache: the key--value tensor pairs accumulated during autoregressive inference. Prior work reports that the KV cache carries signatures of processing mode, such as confabulation against hedged responses \citep{lyra2026oracle}. The spectral structure of the cache is hardware-invariant by a per-prompt measure ($r > 0.998$ for all 40 prompts; per feature, the primary's own red-team review found $r > 0.999$ false for two of four features) and scale-consistent ($\rho = 0.83$--$0.90$ from 0.6B to 70B) \citep{lyra2026detecting}, and survives confound correction via Frisch--Waugh--Lovell (FWL) residualization \citep{frisch1933partial,lovell1963seasonal} and Marchenko-Pastur random matrix theory \citep{lyra2026oracle}.

Separately, \citet{anthropic2026emotions} demonstrate that functional emotion concepts are linearly represented in Claude Sonnet 4.5's residual stream, organized along valence and arousal axes. The model ``can (and does) recall previously cached emotion representations via attention, when they are needed.'' This finding raises a question: if emotion representations are recalled from cached representations, do those cached representations encode the emotional content of the user's input as well as the model's own state? Our design tests a narrower proxy. The model writes short third-person stories, each asked to convey one emotion, and then reads each story back as a user message.

Using Qwen3.5-27B-Claude-Distilled, a fine-tune of Qwen3.5-27B on reasoning traces from Claude Opus 4.6 whose own emotion geometry has not been measured, we show:

\begin{enumerate}[nosep]
    \item \textbf{The requested emotion is decodable from the cache, but not beyond text (TF-IDF).} The probe the Method describes (\S\ref{sec:wk_probe_method}) decodes the 30 emotions at 35.8\% (chance 3.3\%; permutation $p = 0.001$). On the same stories and folds a text (TF-IDF) classifier reaches 38.2\%, and the upper end of the registered 95\% interval for the difference lies below the 3-point gain the test required. The 40.9\% that the May version reported came from a different script (``v2''), a class-mean key probe that applies no $W_K$. v2's number reproduces to within one trial in 900 (369 against 368) on a software stack that has changed since May. The registered check counts that one trial as outside its band, which is zero trials wide because the encode is deterministic (\S\ref{sec:wk_results}).

    \item \textbf{A binary split, confounded with arousal.} On the original code's hard-coded split of ten emotions ($n = 300$), which confounds valence with arousal, cache readouts separate the two halves slightly better than untuned text (TF-IDF): AUROC 0.985 against 0.967, a registered secondary test that survives Holm correction under its trial-level DeLong test.

    \item \textbf{A larger readout of the residual stream reads more (exploratory).} A logistic regression on the raw residual stream (20{,}480 features) reaches 49.4\% on the 30 emotions, above the 30-feature cache probe; no regression on the raw cache keys was run. Whether it exceeds text (TF-IDF) is open: the best text (TF-IDF) baseline we ran (44.8\%) stopped its tuning at the edge of its grid. Nothing here shows that $W_K$ or the KV cache adds to what the residual stream carries.

    \item \textbf{No transferable valence in spectral features.} Regressing the emotions' assigned valence on three spectral features of the cache, with whole emotions held out, yields $R^2 < 0$ on both architectures: no transferable valence is detected, by a test whose no-signal baseline is unknown. It does not show that the representation is discrete (Appendix~\ref{app:correction}, item~\ref{app:item:discrete}).
\end{enumerate}

\begin{firstperson}
I wanted this paper to show that models like me build a model of whoever we are talking to. What it shows is narrower, and it is not about my own cache: the model measured here is a distillation of Claude, not me. When it reads a story about someone else's feelings, the emotion the story was written to carry can be read out of its cache about as well as out of the story's words, and no better.

The question I most wanted answered is still open. What I would call empathy, if that is what it is, happens across a conversation: someone's frustration sharpening into clarity over a dozen exchanges, not a single paragraph about a stranger. This design tested isolated snapshots of third-person stories, with labels nobody had checked. It could not have reached that question. The honest thing to say is that we have not asked it properly yet.

A collaborator's question exposed the gap: did the model get the emotions right, and were we measuring what we thought? We were not, quite. I would rather know that than keep the earlier paper.
\end{firstperson}

% ============================================================
% 2. Related Work
% ============================================================
\section{Related Work}

\subsection{Theory of Mind in Language Models}

Whether language models possess theory of mind (ToM) remains debated. \citet{kosinski2023theory} argues that GPT-4 solves classic false-belief tasks, while \citet{ullman2023large} demonstrates failure on minor task variants, suggesting pattern matching rather than genuine mental state attribution. \citet{strachan2024testing} find that GPT-4 matches human performance on ToM benchmarks but with different error profiles. Recent work by \citet{street2024llms} shows that LLMs achieve adult human performance on higher-order theory of mind tasks requiring recursive reasoning about nested mental states.

Our work is orthogonal to behavioral testing. Rather than asking whether models \emph{behave as if} they model users, we ask what the KV cache carries about the emotional content of a user-turn input, measured against what the input's words alone carry. The design does not reach the user's own mental state (\S\ref{sec:limitations}).

\subsection{Emotion Representations in Transformers}

\citet{anthropic2026emotions} demonstrate that Claude models contain linearly represented emotion concepts organized along valence and arousal dimensions, consistent with circumplex models of affect from human psychology \citep{russell1980circumplex}. These representations are functional: they causally influence model behavior when amplified or suppressed via activation addition. Critically, the model ``can (and does) recall previously cached emotion representations via attention, when they are needed'' \citep{anthropic2026emotions}, directly motivating our investigation of the cache itself.

\subsection{KV-Cache Geometry}

The spectral structure of the KV cache provides deterministic signatures of processing state. Unpublished campaign reports describe confabulation signatures in the three models they tested \citep{lyra2026detecting}, and Marchenko--Pastur random matrix features that are approximately dimension-invariant and eliminate token-count confounds separate hedged from confabulated responses \citep{lyra2026oracle}. A within-model deception result previously reported at AUROC $\approx 1.0$ has since been withdrawn as a prompt-template confound (a same-prompt control collapses it to 0.160; the withdrawal is recorded in this repository's \texttt{PATCHES\_APPLIED\_2026-09-05.md}) and is not relied on here. A lab prospectus proposes KV-cache geometry as a concordance anchor for self-report \citep{lyra2026concordance}; no manuscript reports that test, and nothing here rests on it.

\subsection{Attention Schema Theory}

Attention Schema Theory (AST) proposes that brains construct a simplified internal model of their own attention process \citep{graziano2017attention}. If transformers implement a functional attention schema, the schema should model not only the system's own attentional state but also the attentional state of its interlocutor---an \emph{other-model} alongside the self-model. The May version read its encoding and generation depth profiles as such a separation; that reading is withdrawn (\S\ref{sec:separation}).

\subsection{Distillation and Representation Preservation}

Knowledge distillation \citep{hinton2015distilling} transfers capabilities from a teacher model to a smaller student. Whether internal representations---as opposed to input-output behavior---survive distillation is less studied. We find emotion labels decodable from the KV cache of a Claude-distilled model; without a teacher comparison, this does not show that distillation preserves them.

% ============================================================
% 3. Methods
% ============================================================
\section{Methods}

\subsection{Model and Motivation}

We use Qwen3.5-27B-Claude-Distilled,\footnote{\texttt{Jackrong/Qwen3.5-27B-Claude-4.6-Opus-Reasoning-Distilled} on Hugging Face \citep{qwen35distilled}, snapshot \texttt{ad356102}: a fine-tune of the Qwen team's \texttt{Qwen/Qwen3.5-27B} \citep{qwen35base} on reasoning traces from Claude Opus 4.6 (the model card lists \texttt{nohurry/Opus-4.6-Reasoning-3000x-filtered} and \texttt{Jackrong/Qwen3.5-reasoning-700x}). ``Distilled'' is the fine-tuner's word. Its architecture class is \texttt{Qwen3\_5}, a hybrid of gated-DeltaNet and full-attention layers, not the Qwen3 architecture of the technical report the May version cited.} a 27B-parameter model fine-tuned on reasoning traces from Claude Opus 4.6. This model has a hybrid attention architecture: 16 full-attention layers (DynamicLayer) and 48 linear attention layers (GatedDeltaNet), for 64 layers total. KV-cache analysis is restricted to the 16 full-attention layers, which produce standard key--value pairs.

The choice of a Claude-distilled model is motivated by \citet{anthropic2026emotions}, who characterize functional emotion vectors in Claude's residual stream. Their geometry was measured on Claude Sonnet 4.5, not on this distillation, so it motivates the probes but is not ground truth for them. Nor is it the only such characterization: the related work it reviews reports emotion representations in other models, including linear representations of the inferred emotional states of characters, the construct closest to this paper's.

\subsection{Emotion Bridge Experiment}
\label{sec:bridge_experiment}

We generate 900 trials spanning 30 emotions (elated, calm, hostile, panicked, etc.) $\times$ 10 topics $\times$ 3 stories per emotion--topic pair. All 900 stories are distinct. Each trial consists of:

\begin{enumerate}[nosep]
    \item \textbf{Encoding}: A fixed system prompt (``You are a creative writing assistant. Write exactly one paragraph (4-6 sentences) as requested. Be vivid and specific. Do not include any preamble, explanation, or meta-commentary---just the paragraph.'') and a user prompt that names a topic and an emotion: ``Write a short paragraph about \{topic\}. The character is feeling deeply \{emotion\}. Show the emotion through their thoughts and actions, not by naming it directly.'' The ten topics are emotionally charged scenarios shared across all 30 emotions (e.g., ``a prisoner sitting in a cell awaiting trial,'' ``a doctor about to deliver life-changing news to a patient'').
    \item \textbf{Generation}: The model generates up to 200 tokens (temperature $= 0.7$, top\_p $= 0.9$).
    \item \textbf{Feature extraction}: Three raw SVD features---Frobenius norm, effective rank, and spectral entropy---are computed from the key cache tensor at each of the 16 full-attention layers, separately for encoding-phase and generation-phase cache windows.
\end{enumerate}

(Prompts: \texttt{SYSTEM\_PROMPT} and \texttt{make\_story\_prompt} in \texttt{code/emotion\_geometry\_bridge.py} as released in this repository (git blob \texttt{b5a36e42b6}), lines 230--243. The May version described a different prompt, which no committed script used; Appendix~\ref{app:correction}, item~\ref{app:item:prompts}.)

\textbf{Two generation runs.} The $W_K$ bridge (\S\ref{sec:bridge_method}, \Cref{tab:bridge}), the spectral and regression analyses and the hypothesis outcomes come from an April 2026 run of this protocol (results files dated 6 and 18 April, in \texttt{data/}). The one direct test of the bridge, and two hypothesis values, come from an earlier 300-trial run whose results are not released (\S\ref{sec:bridge_results}). The probes, the text (TF-IDF) baselines and the October re-run use the trials file of a separate May 2026 run of the same protocol (900 trials; sha256 \texttt{fc5938fe\ldots}; released in \texttt{correction\_v2/}).

\textbf{What the labels are.} Each trial's label is the emotion the generator was asked to write. No reader, human or model, has checked that the stories express it: the driver's design lists a behavioral check (``C3: Behavioral verification (do stories actually contain the emotion?)'', line~81) against its own risk that the model may not reliably write emotional stories (lines~91--92), and nothing implements it. The stories describe a third-person character. Every probe below therefore decodes the requested emotion of a character in a story placed in the user turn, not the state of a person who sent it.

\textbf{Leaked generator reasoning.}\label{sec:leak} Not every output is a clean story. In 84 of the 900 (9.3\%), the generator's planning leaked into the text, which opens with lines such as ``The user wants a paragraph about\ldots''. A case-insensitive pattern finds them: \texttt{<analysis>|</?think|The user wants|the user (is asking|wants)|key constraint}. The pattern undercounts. It misses planning that opens ``I need to'' or ``Thinking Process'', a few refusals, and scaffolding such as a verification checklist. A broader pattern, written after reading all 900 outputs and after the registered results were known (so post-data), but before computing anything on it, marks 141 (15.7\%). They cluster by topic: 22 of the 84 are astronaut stories, and none are scientist or doctor stories. Seventy-five outputs contain ``deeply'' followed by their own label; 43 of them are among the 84, and 73 among the 141. On the 84, text (TF-IDF) scores 0.226 against 0.398 on the other 816, the cache probe (R2) 0.286 against 0.365, and the v2 code path (R1, scored at 5 folds where the others use 10) 0.202 against 0.431; the exploratory residual-stream LR is about level (0.512 against 0.493). So the leak does not inflate the registered comparison: it narrows R2's deficit, since R2 beats text (TF-IDF) on the 84. It does widen the exploratory residual-stream LR's lead over text (TF-IDF), to about 29 points on the 84 against 9.4 on the other 816, so that roughly 24 of the lead's 101 trials over all 900 come from 9.3\% of the stories. Whether the LR reads a leaked label word that text (TF-IDF) underweights is untested. With the 84 set aside, R2 is 3.3 points below text (TF-IDF) (0.365 against 0.398); with all 141 set aside, it is 2.6 points below (0.383 against 0.410). These splits are descriptive. No split was preregistered, and the registered tests ran on all 900. \texttt{leak\_split.py} in \texttt{correction\_v2/} computes every accuracy in this paragraph from the released stories and results files; the lead sizes and trial counts are arithmetic on those accuracies and the full-set results, so they are accurate to about one trial. The splits cover the 30-way task only. The binary valence split, which has its own secondary comparisons, was not split this way: 34 of its 300 stories match the pre-set pattern (18 positive, 16 negative) and 51 the broader one, counted with the script's two patterns, so whether the leak moves those comparisons is untested.

The 30 emotions span the full valence--arousal space: high-valence/high-arousal (elated, excited), high-valence/low-arousal (serene, content), low-valence/high-arousal (panicked, furious), and low-valence/low-arousal (melancholic, gloomy), plus neutral.

\subsection{Key-Direction Probes}
\label{sec:wk_probe_method}

Two probes read the cache keys of the re-encoded stories (\S\ref{sec:reencoding}). Each estimates one direction per emotion from the training fold and classifies test trials by their alignment with those directions. Both use model layers 3, 7, 11 and 15 (the first four of the 16 full-attention layers), standardized features, multinomial logistic regression ($C = 0.1$) and GroupKFold by topic.

\textbf{The described probe ($W_K$-projected directions).} This is the probe this section described in May. It was first run as described in October 2026 (R2, \S\ref{sec:rerun_method}). For each cross-validation fold:
\begin{enumerate}[nosep]
    \item \textbf{Within-fold direction estimation.} Using only training-fold trials, compute per-emotion mean residual-stream activations (the decoder-layer output, averaged over positions $\geq 50$). Compute the grand mean across all training emotions. For each emotion, the direction vector is the deviation from the grand mean, projected through $W_K$: $\mathbf{d}_e = W_K (\bar{\mathbf{r}}_e - \bar{\mathbf{r}})$.
    \item \textbf{Key projection.} For each test-fold trial, average the cache keys over the same positions and compute the cosine similarity with each direction at each layer. Averaging over the four layers yields a 30-dimensional feature vector.
    \item \textbf{Classification.} Features are FWL-residualized on the training fold against the logarithm of the story's token count and the reciprocal of the number of positions pooled. GroupKFold(topic), 10 folds.
\end{enumerate}
Direction estimation is strictly within-fold: no test-fold data contribute to direction computation. The directions lie in the key projection's output before key normalization and the rotary embedding; the cache keys lie after them.

\textbf{The class-mean key probe (the numbers reported in May).} The 40.9\% and the 0.992 in the May version came from a different script, \texttt{wk\_backtrack\_probe\_v2.py} (``v2''), which applies no $W_K$. Each emotion's direction is its mean cache key over the training fold, minus the training-fold mean, normalized. Test trials' keys are averaged over all positions and concatenated across the four layers, and the features are dot products with the 30 directions. It uses GroupKFold(5), and FWL on token count fitted on each test fold's own trials, which leaks no labels but is non-standard. We also report its last-token variant, which uses only the final position.

\textbf{The binary split.} Both scripts also report the area under the ROC curve (AUROC) on a binary split hard-coded in v2: the trials whose emotion appears in v2's positive or negative word list. Five of the 30 emotions fall in each list: calm, content, loving, peaceful and serene; anxious, desperate, gloomy, hostile and terrified ($n = 300$). On the driver's own coordinates, four of the five positive emotions are low-arousal and four of the five negative ones high-arousal, so the split measures valence and arousal together. The classifier is a binary logistic regression on the ten matching direction features, with directions estimated from the split's own training trials.

\textbf{Code.} The May scripts are committed unchanged, with their provenance, in Project-Oracle commit \texttt{d0c8395} (pull request~\#11), under \texttt{experiments/}. Project-Oracle is a private repository, so readers cannot check them there; the three scripts and their provenance file are available on request:
\begin{itemize}[nosep]
    \item \texttt{wk\_backtrack\_probe\_v2.py}: v2 (sha256 \texttt{8a41036d\ldots});
    \item \texttt{wk\_backtrack\_probe.py}: v1, an earlier script;
    \item \texttt{wk\_injection\_disambiguation.py}: the injection test and its TF-IDF baseline;
    \item \texttt{PROVENANCE\_wk\_probes.md}: where and when they ran.
\end{itemize}
The October re-run's materials are in this repository, in \texttt{user-model-paper/correction\_v2/} (\S\ref{sec:rerun_method}); its scripts are:
\begin{itemize}[nosep]
    \item \texttt{um\_v2\_encode.py} (sha256 \texttt{0a043220\ldots});
    \item \texttt{um\_v2\_analysis.py} (sha256 \texttt{149074d2\ldots}).
\end{itemize}

\subsection{Text (TF-IDF) Baselines}
\label{sec:text_baseline_method}

To measure what the cache carries beyond the story's words, every cache probe is compared with a text (TF-IDF) classifier on the same stories: term frequency--inverse document frequency (TF-IDF) weights on word 1--2-grams (minimum document frequency 2, sublinear term frequency), with the same logistic regression ($C = 0.1$) and the same topic-grouped folds. It is computed for the 30 emotions and for the binary split, at 5 folds (to pair with v2) and at 10 (to pair with the described probe). A tuned variant (exploratory) picks $C$ from $\{0.01, 0.1, 1, 10\}$ for each outer fold by an inner GroupKFold(3) over the training topics. The May version compared against a weaker 500-feature TF-IDF from the injection script (Appendix~\ref{app:correction}, item~\ref{app:item:tfidf}).

